\documentclass[12pt, a4paper]{article}
\usepackage[utf8]{inputenc}
\usepackage[T1]{fontenc}
\usepackage[polish]{babel}
\usepackage{amsmath}
\usepackage{amssymb}
\usepackage{geometry}
\geometry{margin=2.5cm}
\usepackage{graphicx}
\usepackage{float}

\title{Algorytm Dijkstry zrównoleglony metodą delta-stepping}
\author{Karol Helm, Maciej Anioł}
\date{Listopad 2026}

\begin{document}

\maketitle

\section{Wprowadzenie}
Celem projektu było zimplementowanie oraz przebadanie uwspółbieżnionej wersji algorytmu Dijkstry wykorzystującej metodę delta-stepping opisaną w artykule U. Meyer and P. Sanders, $\Delta$-stepping: a parallelizable shortest path algorithm, Journal of Algorithms 49 (2003). 
\\
Dany jest graf skierowany o nieujemnych wagach całkowitych. Obliczamy długości najkrótszych ścieżek z wierzchołka $0$ do wszystkich wierzchołków. Zwracamy odległość do wierzchołka $n-1$ lub $-1$, jeśli wierzchołek $n-1$ jest nieosiągalny. 
\\
Paczka projektu zawiera:
\begin{itemize}
    \item Rozwiązanie jednowątkowe $dijkstraOneThread.cpp$
    \item Rozwiązanie wielowątkowe $dijkstraDeltaStepping.cpp$
    \item Raport $raport.pdf$
    \item Skrypt $run.sh$ uruchamiający $scaling.py$
    \item Program w Pythonie generujący grafy oraz porównujący czasy działania rozwiązań $scaling.py$ 
    \item Prezentację na zajęcia $prezentacja.pdf$
\end{itemize}
\section{Opis wersji jednowątkowej i wielowątkowej}
\subsection{Wersja jednowątkowa}
Wersja jednowątkowa jest implementacją klasycznego algorytmu Dijkstry korzystającego z kolejki priorytetowej. 
\begin{enumerate}
    \item Zdejmujemy z kolejki wierzchołek o najmniejszej odległości.
    \item Sprawdzamy, czy ta odległość jest aktualna.
    \item Relaksujemy krawędzie i jeśli poprawiamy odległość do innego wierzchołka to wrzucamy go na kolejkę.
\end{enumerate}
Konieczność utrzymania konkretnej kolejności przetwarzania wierzchołków utrudnia efektywne zrównoleglenie tej wersji algorytmu.
\subsection{Wersja wielowątkowa}
Wersja wielowątkowa korzysta z techniki delta-stepping. Zamiast wybierania pojedynczego wierzchołka o minimalnej odległości, algorytm przetwarza wierzchołki według kubełków odpowiadających przedziałom odległości
$[i\Delta, (i+1)\Delta)$. Współczynnik $\Delta$  może być dowolną dodatnią liczbą, ale jego wybór ma wpływ na czas działania algorytmu. W naszej implementacji to $max(1, floor(2W/m))$, gdzie $W$ to suma wag, a $m$ to liczba krawędzi. Dla $m=0$ przyjmujemy $\Delta=1$. Przed rozpoczęciem właściwego algorytmu dla każdego wierzchołka dzielimy krawędzie na lekkie ($waga <= \Delta$) oraz ciężkie ($waga > \Delta$). Lekkie krawędzie mogą poprawiać odległości
wewnątrz aktualnego kubełka, natomiast ciężkie na pewno prowadzą
do późniejszych kubełków od aktualnie przetwarzanego. Pozwala to odroczyć ciężkie relaksacje do zakończenia fazy lekkiej i ograniczyć liczbę operacji. Zaczynamy od jednego kubełka zawierającego wierzchołek $0$ z odległością $0$. Następnie dopóki wszystkie kubełki nie są puste:
\begin{enumerate}
    \item Wybieramy najmniejszy niepusty kubełek. 
    \item Relaksujemy równolegle krawędzie lekkie zapisując ich wynikowe pozycje do bufora
    \item Jeden wątek przenosi pozycje z buforów do kubełków. Jeśli aktualny kubełek nie jest pusty, to powtarzamy proces relaksacji lekkich krawędzi
    \item Po zakończeniu fazy lekkiej relaksujemy równolegle krawędzie ciężkie. Korzystamy z tablicy wartości $added$, aby każdy wierzchołek trafił do fazy ciężkiej tylko raz
    \item Ponownie jeden wątek przenosi wierzchołki z buforów do kubełków
\end{enumerate}
Wersja wielowątkowa korzysta z biblioteki OpenMP. Przetwarzanie kubełków odbywa się w bloku \texttt{omp parallel}. Pętla, która przechodzi przez wierzchołki obecnego kubełka, jest podzielona przez \texttt{omp for}. Bufor przy relaksacji jest konieczny, ponieważ operacja \texttt{push\_back} nie jest atomowa. Przy scalaniu buforów korzystamy z \texttt{omp single}. Odległości przechowujemy w wektorze elementów atomowych. Funkcja \texttt{compare\_exchange\_weak} pozwala bezpiecznie zaktualizować odległość, kiedy kilka wątków jednocześnie
znajduje różne drogi do tego samego wierzchołka. Sekwencyjność pozostaje przy wczytaniu i inicjalizacji danych, budowie grafu, podziale na krawędzie light i heavy oraz wcześniej wspomnianym scalaniu buforów. 
\section{Pomiar wydajności algorytmu}
Wydajność była badana zarówno na urządzeniu lokalnym (CPU: Intel Core i7-1260P), jak i na serwerze (CPU: 4 x Xeon Gold 6154). Wyniki prezentują się następująco:
\begin{figure}[H]
    \centering
    \includegraphics[width=\textwidth]{wyniki2.png}
    \caption{Mediany czasu działania algorytmów na serwerze i lokalnie,
    bez operacji wejścia oraz wyjścia. Czasy podano w sekundach. Tabelę przygotowano w programie Excel.}
    \label{fig:wyniki-serwer}
\end{figure}
Rozwiązania zostały przebadane na: 
\begin{enumerate}
    \item Trzech grafach Random - losowe grafy na 4 mln wierzchołków i na 16 mln krawędzi. Rozwiązania zostały uruchomione na każdym grafie z tej grupy po trzy razy, po czym obliczono medianę.
    \item Grafie Path - ścieżka na 10 tys. wierzchołków, zostało wykonanych po 5 pomiarów i wybrano medianę.
    \item Grafie warstwowym na 2 mln wierzchołków. Wykonano 5 pomiarów i wpisano do tabeli medianę.
\end{enumerate}
Analiza uzyskanych danych znajduje się w sekcji Wnioski.
\section{Kompilacja i uruchomienie}
Domyślne uruchomienie programu następuje poprzez wywołanie skryptu $run.sh$, który uruchamia $scaling.py$ - program w Pythonie odpowiedzialny za skompilowanie rozwiązań oraz wygenerowanie grafów testowych, na których następnie czasy rozwiązań są porównywane. Oczywiście można też skompilować i przetestować rozwiązania ręcznie komendami: \begin{verbatim}
g++ src/dijkstraOneThread.cpp -O3 -std=c++17 -o dijkstraOneThread
g++ src/dijkstraDeltaStepping.cpp -O3 -std=c++17 -fopenmp -o dijkstraDeltaStepping
\end{verbatim}
Oba programy przyjmują wejście w postaci kolejnych linii:
\begin{enumerate}
    \item liczba zestawów danych $z$
    \item liczba wierzchołków $n$ i krawędzi $m$
    \item $m$ linii opisujących krawędzie skierowane w postaci $u$, $v$, $val$.
\end{enumerate}
\section{Wnioski}
Nasza implementacja zrównoleglonego algorytmu Dijkstry przy uruchomieniu na serwerze uzyskała przyspieszenie $3{,}61\times$, $3{,}29\times$, $3{,}31\times$ na grafach losowych oraz $2{,}22\times$ na grafie warstwowym. Przy uruchomieniu lokalnym uzyskano przyspieszenie $1{,}61\times$, $2{,}82\times$, $2{,}84\times$ na grafach losowych
oraz $2{,}59\times$ na grafie warstwowym. Mnożnik został wybrany jako najbardziej efektywne dobranie liczby wątków przy każdym grafie. Niestety delta-stepping okazuje się nie być efektywny dla ścieżki. Straty wynikikające z wielowątkowości znacznie przewyższają zyski, przez co program dla ścieżki działa wolniej.
\end{document}
